\chapter{TIC TOC, TIC TOC}\label{tic-toc-tic-toc}

\hfill\break

«C'è qualcosa che posso fare per te, sergente Bishop?», mi chiese
Kendall il pomeriggio successivo.

«Oh. Ehm, niente. È un onore servire, sergente scelto.»

Lei inclinò la testa. Era chiaro che non si era bevuta le mie stronzate:
«Tutti noi serviamo, a modo nostro. Tu hai salvato il mondo. Ho parlato
con i miei genitori al telefono ieri sera, e mio padre è scoppiato a
piangere, era così sollevato nel sapere che non siamo più sotto il
controllo delle lucertole. Mio padre è l'uomo più duro che conosca; era
un ranger dell'esercito, ha perso una gamba dal ginocchio in giù in
Afghanistan. Dissero che non sarebbe mai più stato in grado di camminare
bene, lui fece buon viso a cattivo gioco e si qualificò per il servizio
di fanteria diciotto mesi dopo. Si è qualificato e ha servito, con mezza
gamba. L'uomo più duro che conosca, e piangeva, parlando con me. Mi ha
detto che pensava che noi umani non saremmo mai stati altro che schiavi,
se fossimo sopravvissuti. Ora abbiamo di nuovo speranza e non è dovuto a
nulla di quello che abbiamo fatto noi quaggiù. È dovuto al miracolo
della tua apparizione in orbita, per vaporizzare quelle lucertole figlie
di puttana, in qualunque modo tu l'abbia fatto. Quindi», fece un respiro
profondo e mi guardò con piglio severo, «sergente, c'è qualcosa che
posso fare per te?»

Capii. Non si trattava di me. Mi ricordai di qualcosa che avevo sentito
da un tizio che aveva ricevuto la Medaglia d'onore: indossare la
medaglia non è qualcosa che si fa per se stessi. Mette in imbarazzo le
persone attorno a te e ti mette a disagio, e alza una barriera tra il
premiato e tutti gli altri.

E non vale solo per la Medaglia d'onore, vale per qualsiasi medaglia
assegnata al valore in combattimento. Non indossi una medaglia per te
stesso. La indossi per i ragazzi che non ce l'hanno fatta. La indossi
per la tua unità, per il tuo servizio, per il tuo paese. Non si tratta
di te, si tratta delle persone che ti assegnano la medaglia. Si tratta
del loro bisogno di esprimere gratitudine per le tue azioni in modo
tangibile. Il sergente scelto Kendall aveva bisogno di un modo per
sentire che aveva fatto qualcosa, qualsiasi cosa, per me, per ricambiare
qualunque cosa avessi fatto.

Non si trattava di me, non importava quanto mi mettessero a disagio
l'attenzione e le cerimonie. Non volevo una medaglia, e non\ldots{}
Un'idea mi colpì come una padella di ghisa sbattuta sulla testa: «Vorrei
un cheeseburger», ammisi. «Non so dirle quanto vorrei un cheeseburger.
Non mangio un cazzo di cheeseburger da quando sono salito sull'ascensore
spaziale e ho lasciato la Terra. Per tutto il tempo che siamo stati in
transito, su Campo Alfa, su Paradiso, neanche uno straccio di
cheeseburger in vista. Sto parlando di un vero e proprio cheeseburger
americano, non di una roba da fast food. Un cheeseburger cucinato su una
griglia a carbone nel cortile sul retro il 4 luglio.» Mi resi conto che
stavo divagando, sbavando, ma non potevo farne a meno. «Un patty di
manzo home made, non troppo grande, non troppo spesso, senza
schiacciarlo al punto da renderlo denso come un disco da hockey.
Grigliato al punto giusto, come se fosse al sangue, poi ci si mette
sopra una fetta di cheddar e la si fa sciogliere un po', così il
formaggio fa le bollicine ma non si fonde del tutto. Il panino rotondo
non dev'essere spesso, non uno di quei rotoli Kaiser o robe tipo
brioche, il panino serve solo per tenere tutto insieme, non è la star
dello spettacolo. Grigli un poco il panino, ma non come per fare del
pane tostato: solo per renderlo un po' croccante. E aggiungi cipolle
alla griglia e ketchup. È tutto quello che serve a un buon
cheeseburger.»

Kendall mi lanciò uno sguardo che non seppi interpretare. Forse mi ero
lasciato trasportare dall'entusiasmo per gli hamburger. Poi annuì,
sorrise, e giuro che le si piegarono appena un po' le ginocchia. «Oh, so
esattamente quello che vuoi dire. Adoro un buon cheeseburger. Stasera
saltiamo la cena in sala mensa e tu vieni nei nostri alloggi, siamo
dietro la casa del comandante. Abbiamo una griglia, e ti preparo un vero
cheeseburger.» Una delle guardie che erano con lei si schiarì la gola e
Kendall gli lanciò un'occhiata. «Sì, di questi tempi è difficile trovare
del vero manzo, ma per te, sergente, l'Aeronautica degli Stati Uniti
farà un'eccezione, questo è sicuro. È il minimo che possiamo fare.»

\hfill\break

Fedele alla sua parola, il sergente Kendall mi portò nell'edificio in
cui alloggiava, e sul retro c'era una griglia. Faceva un po' freddo
nella notte ad alta quota di Colorado Springs, il che significa solo che
dovevamo indossare delle giacche. Nessuno di noi avrebbe perso
l'occasione di un barbecue. Tenni il cheeseburger con entrambe le mani e
inalai a fondo. Era quasi perfetto. Se fosse stato cucinato da mio padre
su una griglia nel cortile dei miei genitori, sarebbe stato il
cheeseburger migliore di tutti i tempi, perfetto al 100\%, ma un legame
familiare era l'unica cosa che mancava. Diedi un morso esitante: «Oooh.
Oooh, bello, quanto è buono. Non ha un'idea di quanto abbia sognato
questo momento».

Kendall sollevò il suo cheeseburger alla salute: «Anche noi era da un
po' che non ne mangiavamo, quindi grazie per questo». I componenti della
sua squadra annuirono, mentre masticavano estasiati.

Avevo quasi finito di mangiare quando mi accigliai.

«Che cosa c'è?», chiese Kendall.

«Beh, stavo pensando che i componenti della Expeditionary Force su
Paradiso potrebbero non mangiare più cheeseburger. Mai più. Dovrebbero
essere in grado di coltivare abbastanza cibo, ma saranno vegetariani in
senso stretto. Anche prima che i ruhar riconquistassero il posto, le
spedizioni di rifornimenti dalla Terra si erano fermate.» Questo
significava niente medicine. Speravo che i ruhar condividessero la loro
tecnologia medica avanzata più di quanto non avessero fatto i kristang.

«Sì», Kendall annuì tristemente, «le lucertole non ci hanno detto nulla,
ma l'ascensore spaziale ha smesso di funzionare un paio di mesi dopo che
eravate partiti, era ovvio che non vi spedivano provviste. Finché non
siete tornati, non avevamo idea di cosa stesse succedendo nell'Unef. Le
maledette lucertole non ci hanno detto niente. Mio fratello è in
servizio con l'Unef, fa parte del III fanteria. Pensi che staranno
bene?»

Pensai alla Bürgermeister e alle sue promesse di prendersi cura degli
umani su Paradiso. «Non so se sia tutto a posto, ma, sì, penso che
staranno bene con i criceti al comando. Suo fratello farà meglio ad
abituarsi a essere un agricoltore. Ehi», sollevai l'ultima metà del mio
cheeseburger, «all'Unef. Questo lo mangio per loro.»

«Evviva!», gridarono tutti, e ci godemmo in silenzio i nostri hamburger.

«Cazzo, è stato bello», disse Kendall. «Un altro?», guardò verso la
griglia.

«Oh sì», sorrisi. Salvare il mondo ha i suoi vantaggi.

\hfill\break

«Alzati e splendi, dormiglione!», annunciò Skippy, svegliandomi da un
sonno profondo la mattina dopo.

«Uffa. Che ora è?», chiesi assonnato.

«Quindici minuti prima di quando il sergente scelto Kendall ha
intenzione di svegliarti.»

«Allora ho tempo per altri dieci minuti di riposo», mi misi il cuscino
in testa per levarmelo dai piedi.

«Non esiste, Joe, questo è il nostro momento. Ti vanno le coccole?»

«Neanche se tu fossi una vera lattina di birra.»

«Mi ferisci. Ehi, vuoi sapere cosa ho fatto ieri notte, mentre dormivi?»

Oh merda. Mi tirai su a sedere e gettai il cuscino dall'altra parte
della stanza. «Che cazzo hai combinato stavolta, Skippy? Ti annoiavi e
sei entrato nei file secretati delle agenzie di altri governi?»

«Eh? No, l'ho fatto l'altra sera, mentre leggevo i file dell'Agenzia per
la sicurezza nazionale. Non preoccuparti, nessuno degli altri governi
del mondo ha segreti che valga la pena mantenere.»

«Ma non eri impegnato a chiacchierare con la gente, tipo con tutta la
gente?»

«Sì, grazie, lo faccio ancora, e la maggior parte di loro pensa ancora
che sia uno stronzo.»

«Non l'avrei mai detto. Adesso hai un campione abbastanza ampio per
determinare che sei davvero uno stronzo?»

«La giuria è ancora fuori per deliberare. Considerando che si tratta di
una giuria piena di scimmie, ho intenzione di fregarmene. Comunque, per
quanto affascinante sia chiacchierare con diversi miliardi di scimmie,
mi annoiavo di nuovo e ho fatto qualcosa di utile. Ho finito di
scaricare tutti i dati.»

«Cosa intendi con ``tutti i dati''?»

«Tutti i dati memorizzati in banche dati accessibili sulla Terra. Ma tu
pensa. Sono solo un paio di exabyte, posso memorizzarli in un'unghia,
per così dire.»

«Porca vacca.» Non avevo idea di cosa fosse un exabyte. Sembrava
impressionante.

«Già. Anch'io ho fatto qualcosa di utile, dopo essermi stancato di
correggere gli incredibili errori logici nelle vostre cosiddette riviste
scientifiche. Ho risolto dei crimini.»

«Adesso sei Sherlock Holmes?»

«Dato che Holmes era molto più intelligente dell'umano medio, sì, perché
no? Ho confrontato le impronte lasciate dai criminali sul luogo del
delitto con quelle lasciate sugli schermi dei loro dispositivi mobili, o
nelle loro case, nel raggio delle webcam. Ho anche confrontato il Dna
nelle banche dati della polizia, è scioccante quanto siano pessime le
vostre forze dell'ordine nella condivisione dei dati. In alcune
occasioni, mi è bastato leggere gli appunti dei casi per capire chi
fosse il colpevole, usando tutti i dati disponibili, che ora sono tutti
quelli memorizzati in formato elettronico. Se i vostri poliziotti
alzassero i loro pigri culi e analizzassero i kit di Dna già in loro
possesso, potrei risolvere molti più crimini.»

Avevo letto da qualche parte quanti campioni di Dna, compresi quelli
raccolti dai kit stupro, non fossero stati testati. «È una questione di
risorse, Skippy, la polizia è\ldots»

«No, è una questione di priorità. La vostra specie non pensa che
ottenere giustizia per le vittime di un crimine sia abbastanza
importante da finanziare in modo adeguato i laboratori forensi. E alcuni
poliziotti in tutto il mondo sono corrotti al 100\%. Voi umani non mi
impressionate in questo campo.»

«Non posso contraddirti su questo, Skippy.»

«Ehm.» Sembrava deluso di perdere l'opportunità di una discussione.
«Quindi, ecco il problema: ho risolto oltre sessantamila crimini e ora
non so cosa fare con i dati ottenuti, non potendomi rivelare al resto
della Terra.»

«Ah. Ehm, puoi mandare tutto all'Fbi? Là ci dev'essere qualcuno
autorizzato a sapere di te.» Ero confuso su chi sapeva cosa. Il sergente
scelto Kendall e la sua squadra di sicurezza mi seguivano ovunque, ed
era chiaro che sapevano di Skippy, ma non sapevo quanto sapessero o
dovessero sapere.

«L'arretrato di crimini irrisolti dell'Fbi è parte del problema. E
questo non riguarda altri paesi.»

Non avevo idea di cosa fare, non avevo esperienza in applicazione della
legge. O sicurezza delle informazioni. «Skippy, ne parlerò con qualcuno.
Penso sia fantastico che tu stia usando le tue, ehm, risorse, il tuo
talento, sai, per aiutare. Aiutare la gente ad avere giustizia.»

«Ancora più importante, togliere i criminali dalle vostre strade. Ci
sono un sacco di recidivi là fuori, che non sono mai stati arrestati.»

«Lo saranno. Ehi, ora che hai risucchiato tutti i dati del pianeta e
risolto migliaia di crimini, cos'hai intenzione di fare per tenerti
occupato? Chiacchierare con miliardi di umani non ti basta, vero?»

«Neanche lontanamente. È divertente e abbastanza interessante per ora,
mi impedisce di sentirmi solo. Ti ringrazio per questo, è stata una
buona idea.»

«Ti impedisce di sentirti solo, ma il resto di te si annoia?»

«No, il resto di me è in standby. Quello che tu pensi come Skippy è una
minuscola, minuscola sotto-mente che ho creato per gestire le nostre
interazioni. È in questo modo che sono sopravvissuto da solo per così
tanto tempo senza impazzire; ho creato una sotto-mente per controllare
ogni tanto se qualcosa fosse cambiato, e il resto di me in pratica ha
dormito, a lungo. Mi sono emozionato quando la mia sotto-mente ha
rilevato la prima nave kristang che saltava nel sistema di Paradiso, e
di nuovo quando le lucertole mi hanno estratto dal terreno. Poi quegli
idioti mi hanno messo su uno scaffale in un magazzino, e quando i ruhar
hanno conquistato il pianeta mi hanno dato un'occhiata e mi hanno
rimesso sullo scaffale. Sono stato attivo con continuità solo da quando
i primi umani sono sbarcati su Paradiso.»

«Questa tua sotto-mente gestisce tutto?» Immaginai un ragazzo su un
divano, che guarda una partita di calcio, naviga su internet, messaggia
con gli amici e poi, distrattamente, parla con la moglie bevendo una
birra. Era Skippy che parlava con me, anche se, a differenza di un
ragazzo distratto, Skippy poteva interagire pienamente quando mi
serviva. «Prendere il controllo della nave thuranin, programmare salti,
deformare lo spazio-tempo, tutto questo?»

«No, impiego altre risorse al bisogno, non ho mai usato più del 7\%
delle mie capacità da quando ti ho incontrato. I miei archivi indicano
che il massimo che abbia mai usato è il 62\%, il che m'induce a
chiedermi quanto siano accurati i miei archivi e per quale scopo sia
stata progettata la mia capacità. Perché necessito di tutta questa
memoria ed enorme potenza di elaborazione? Joe, questo è il motivo per
cui devo contattare il Collettivo. Ho bisogno di sapere chi sono.»

\hfill\break

Le due settimane successive furono un caos. Trascorsi la maggior parte
di altri tre giorni impegnato in interrogatori divenuti così ripetitivi
che anche gli informatori avevano finito le nuove domande da fare.
Skippy aveva dato alla Cia un enorme dump di informazioni da
distribuire, così avrebbero smesso di fargli domande stupide. Con
l'accesso al dump dei dati, qualsiasi cosa memorizzata nel mio cervello
sarebbe stata ridondante. Dopo di che, volammo a Parigi a bordo dell'Air
Force One. I maggiori leader mondiali si riunivano lì per una
conferenza, per discutere il da farsi.

All'incontro portai anche Skippy, che naturalmente chiacchierava con i
vari leader nella lingua madre di ciascuno. Poi questi ultimi si
ritirarono per parlare a porte chiuse. E parlare. E parlare. Avevano
molto di cui parlare.

«Tic toc», mi diceva Skippy ogni mattina, «tic toc.» L'orologio stava
ticchettando, misurando il tempo che mancava al riavvio del wormhole
locale. Il terzo pomeriggio della conferenza, scoprii che Skippy aveva
monopolizzato lo schermo video principale con la scritta ``L'OROLOGIO
STA TICCHETTANDO, SCIMMIE!'' in diverse lingue.

Ricevettero il messaggio.

\hfill\break

I governi possono metterci una vita a prendere una decisione, anche
quando è chiaro come il sole che cosa bisogna fare. Una volta presa la
decisione, però, le cose possono muoversi in fretta. Il generale Brenner
mi chiamò nel suo ufficio temporaneo a Parigi. Sembrava occupato, con
una serie di alti ufficiali che entravano e uscivano dal suo ufficio, ma
furono tutti cacciati via quando arrivai con Skippy: «Bishop, andrò
dritto al punto. Stiamo per far decollare l'Olandese con un equipaggio
internazionale», un'espressione aspra gli balenò in faccia, «devo sapere
se ti unirai a loro».

«Cosa? Signore, voglio dire, sì, assolutamente. Ho promesso a Skippy che
avrei trovato questo Collettivo. Ho fatto un accordo e gli terremo
fede.» Se fossi stato autorizzato a fare un accordo del genere era una
cazzo di domanda inutile, di cui avrebbero potuto discutere i
quarterback del lunedì mattina. «Ha intenzione di comandare la missione,
signore?» Pensavo che, equipaggio internazionale o no, al comando ci
sarebbe stato un americano. Brenner o un generale dell'Aeronautica o un
ammiraglio della Marina. L'esercito non aveva mai comandato una nave
madre, ma lo stesso valeva per gli altri servizi armati. Io, ovvio,
tifavo perché fosse l'esercito a ottenere l'onore. Hooah.

«No, Bishop, non sarò a bordo, ho abbastanza merda di cui occuparmi
quaggiù, devo sistemare il casino e assicurarmi che nessuno approfitti
del caos per alzare la cresta. Sarai tu al comando.»

«Signore?»

Brenner allungò la mano in un cassetto, tirò fuori una piccola scatola
di cartone e la fece scivolare sul tavolo verso di me. «Rimettiti queste
aquile, ripristineremo la tua promozione sul campo a colonnello, per
questa missione. Non montarti la testa.»

«Signore?» Sì, sembravo stupido, ma lo sareste sembrati anche voi se
foste stati sorpresi come lo ero io. La mia promozione su Paradiso era
stata una trovata pubblicitaria a beneficio dei kristang, e tutti lo
sapevano. Adesso era il capo di stato maggiore dell'esercito a dirmi che
ero di nuovo un colonnello, per davvero questa volta.

«Abituati, colonnello Bishop. Gli stati maggiori riuniti ne hanno
discusso con la presidentessa. Se mettessimo qualcun altro al comando,
diventerebbe una lotta per il prestigio nazionale, e noi non abbiamo
tempo per questo. A essere franco, la tua mancanza di esperienza non
significa un bel niente, perché nessun altro ha esperienza al comando di
un'astronave.»

«Puoi giurarci, amico!» Skippy parlò per la prima volta. Lo tenevo
dentro uno zaino, cosa che non cambiava un bel niente: anche se fosse
stato in fondo all'oceano avrebbe ascoltato le mie conversazioni.

«E nessun altro ha esperienza con quelle teste di cazzo», disse Brenner
con un'alzata di sopracciglia, mentre i lati della sua bocca si
sollevavano per farmi capire che sapeva anche lui come trattare con
Skippy. Mi diede un pezzo di carta: «Questa è una lista di volontari per
il suo equipaggio».

Diedi un'occhiata alla lista. Non mi sorprendeva che ogni pirata
dell'equipaggio originale, che non fosse rimasto ferito in modo troppo
grave da non essere idoneo al combattimento, si fosse offerto di nuovo
volontario. Chang, che era ancora segnalato come tenente colonnello,
quindi anche la sua promozione sul campo era stata confermata. Desai,
che era stata promossa a maggiore. Il maggiore Simms. Il sergente scelto
Adams. Tutti. L'elenco comprendeva anche un sacco di forze speciali e
una mezza dozzina di piloti molto esperti. Avrei chiarito molto bene che
Desai era il pilota capo, a meno che non fosse stata lei a dirmi il
contrario. Quello che mi sorprendeva era il numero di civili sulla
lista, per lo più scienziati. «Signore, temo che stiamo portando più
gente di quanta ce ne serva, molta di più.»

«In che senso?», chiese Brenner.

«Non abbiamo bisogno di tutte queste persone», tamburellai con le dita
sulla lista dei nomi, «per raggiungere gli obiettivi della missione. Se
qualcosa va storto, questa», tamburellai di nuovo con le dita sulla
lista, «è solo altra gente che non tornerà a casa.»

Brenner mi lanciò uno sguardo duro, quindi non aspettai. Ero stanco, il
che mi rendeva irritabile. Ancora più importante, avevo bisogno di
sapere fino a che punto avrei potuto tirare la corda con l'esercito
prima che lo facessero loro. Dovevo sapere quanto avevano bisogno di me:
«Questo incarico di comando non mi è stato affidato perché l'esercito ha
grande fiducia nelle mie capacità di leadership. Sono qui perché Skippy
vuole me, e perché sono sacrificabile».

«Va bene, colonnello. Dimmi, quali sono per te gli obiettivi della
missione?», chiese il generale Brenner.

Risposi con cautela, consapevole che si trattava di un test. «Vedo tre
obiettivi. Primo, volare verso il wormhole, attraversarlo e chiedere a
Skippy di chiuderlo per sempre dietro di noi. Bloccarlo, così nessuno
tranne Skippy potrà usarlo di nuovo. Secondo, continuare ad assicurarci
che nessun'altra specie possa mai scoprire che gli umani sono coinvolti
nel furto delle navi kristang e thuranin, perché, se la coalizione
maxolhx dovesse mai venire a sapere la verità, la Terra potrebbe essere
nei guai, anche senza che gli alieni abbiano accesso al wormhole locale.
E terzo, far felice Skippy perché stiamo tenendo fede all'accordo e
facendo la nostra parte, aiutandolo a entrare in contatto con questo
Collettivo, se esiste ancora. Ho detto bene? L'ordine di priorità è
corretto?»

«Sembra corretto», disse Brenner con un sorriso stretto.

«Nessuno di questi obiettivi per la missione implica il nostro ritorno
sulla Terra. Lo scenario ottimale per il pianeta è: attraversiamo il
wormhole, Skippy lo chiude dietro di noi e l'Olandese esplode
all'istante in un miliardo di pezzi.» Guardai Brenner, e lui non mi
disse che mi sbagliavo su questo, perciò continuai. «Se per un qualche
miracolo Skippy individuasse il Collettivo e venisse, ehm, trasportato
fino al paradiso dei computer o qualunque cosa lui pensi che accada»,
Skippy aveva mantenuto un frustrante silenzio sui suoi piani relativi a
quello che sarebbe successo dopo che avesse contattato il Collettivo,
«allora ci troveremmo di fronte alla sfida infernale di cercare di
manovrare o pilotare l'Olandese. È probabile che resteremmo bloccati
nello spazio, e dovrei ordinare alla nave di autodistruggersi, per
evitare che un giorno ci catturino. Quindi», guardai di nuovo Brenner
negli occhi, sapendo che aveva preso decisioni e dato ordini che avevano
mandato uomini a morire in combattimento prima, «non voglio portare con
noi più persone di quelle di cui abbiamo assoluto bisogno.»

«Ognuno su quella lista è un volontario, e tu hai bisogno di una forza
considerevole, perché non hai idea di quello che incontrerai là fuori.»

«Volontari che pensano che questa sia una grande avventura, che
salveremo l'umanità e torneremo carichi di conoscenza e tecnologia»,
scossi la testa. «Salveremo l'umanità. È tutto quello che stiamo
facendo. Farò il mio dovere, signore, anche se non tornerò mai più. Non
ha senso chiedere», diedi un'altra occhiata alla lista, «ad altre
settanta persone di correre lo stesso rischio.»

«Bishop», questa volta non c'era accenno di sorriso sulle labbra di
Brenner, «essere al comando significa rischiare la vita delle persone
per raggiungere gli obiettivi della missione. A volte significa mandare
persone buone, dedite, coraggiose, in situazioni da cui è probabile che
non tornino indietro. Se non sei in grado di farlo, non sei l'uomo
giusto per questo lavoro.»

«Signore, l'ho fatto, lo sa.» Due volte. No, tre volte. Quattro in
realtà. A casa, avevo chiesto ai miei vicini di catturare un soldato
alieno, usando fucili e un camion dei gelati. Al Vettore, pensavo
stessimo facendo un gesto inutile, che i criceti ci vedessero nel
condotto di lancio o da qualche parte lungo la strada e ci uccidessero
come bersagli facili. Lo scopo del nostro tentativo di reazione al
Vettore non era realizzare qualcosa di utile contro i ruhar, perché a
quel punto pensavamo tutti che si fossero ripresi il pianeta per sempre.
Il punto era mostrare ai kristang che i loro alleati umani non si
arrendevano nemmeno quando non avevano nessuna chance, in modo che le
lucertole non avrebbero pensato alle persone sulla Terra come a inutili
codardi. ``È probabile che siamo comunque spacciati'', pensavo, ``quindi
perché non contrattaccare?'' Inoltre, ero un soldato. Sebbene mi fossi
arruolato più per pagare l'università che per patriottismo, l'esercito
mi aveva in qualche modo addestrato a essere un soldato, e i soldati non
mollano. La terza volta era stata quando avevamo preso la Fiore: non
avevo creduto davvero che potessimo farlo fino a quando l'ultimo
kristang non era morto e il reattore esploso. E la quarta volta era
stato l'assalto alla base sull'asteroide. Quello era stato il più duro
per me perché, a differenza di quanto avvenuto durante l'azione al
Vettore e la cattura della Fiore, in quell'assalto la mia persona non
era a rischio. Chang e Giraud e Thompson e Adams e gli altri erano a
rischio, li avevo mandati là fuori mentre il mio culo era al sicuro a
bordo dell'Olandese, con le dita di Desai sui comandi; eravamo pronti a
saltare via se qualcosa avesse minacciato la nave. «Ho rischiato la vita
delle persone anche quando pensavo che non avessimo possibilità di
successo. Lo rifarò quando lo riterrò necessario. Quello che non farò è
ingannare la gente. Signore, parlerò con tutte le persone incluse in
questa lista», Skippy avrebbe potuto gestire la traduzione per me, «e
dirò loro in tutta onestà cosa penso di questa missione. Se vorranno
ancora partire, allora sarò onorato di servire con loro.»

Brenner annuì. «Sono sicuro che tutte queste persone sanno quello che
stai per dire, ma vai avanti, non saprei suggerirti altrimenti. Bishop,
la migliore squadra che puoi portare in combattimento è quella composta
da uomini e donne che conoscono i rischi e ti seguono comunque.»

«Quando partiamo?» La mia mente stava correndo attraverso tutte le cose
che dovevano accadere prima che l'Olandese entrasse in orbita. Ed era
solo quello a cui riuscivo a pensare, ero sicuro di stare dimenticando
un milione di cose importanti. Sarebbe stato bello affidarsi a
un'intelligenza artificiale super-brillante, ma Skippy era troppo
distratto. Aveva anche dimostrato molte volte che non gli riusciva
facile pensare come noi rifiuti biologici. Se avessi affidato la
logistica a Skippy, si sarebbe ricordato di tutto tranne che dell'acqua.
O dell'ossigeno.

«Dopodomani. Dobbiamo sbarazzarci di quei due kristang a bordo della
nave da trasporto truppe, colpire i loro tre siti qui con i cannoni a
rotaia, caricare le provviste a bordo dell'Olandese e avere ancora a
disposizione un sacco di tempo perché possiate viaggiare verso il
wormhole prima che si resetti.» Guardò fuori dalla porta aperta
l'assistente, che aveva cercato di attirare la sua attenzione. Scosse la
testa: «E prima che la nostra leadership civile cambi idea in proposito.
Ma ancora prima di volare fino all'Olandese, devi incontrare la
leadership civile delle nazioni coinvolte e il tuo equipaggio
volontario. Il tenente colonnello Chang e altri sono già qui. Il resto
delle persone sta arrivando».

«In tal caso, signore, vorrei chiedere a Skippy di portare qui la
navicella thuranin dal Colorado, così posso volare dritto fino
all'Olandese.»

«Nessun problema, colonnello Joe!», disse Skippy con entusiasmo. «Sto
preparando la navicella in questo momento.»

«D'accordo, colonnello, ma ho un suggerimento», aggiunse Brenner. «Prima
di andare in orbita, fermati in Maine per vedere i tuoi genitori. Non è
ancora stato annunciato, ma i governi coinvolti stanno rinunciando a
cercare di passare sotto silenzio l'accaduto. L'esistenza di Skippy è
ancora un segreto molto stretto e il governo, i governi di tutto il
mondo non riconoscono nulla in via ufficiale. Non possiamo nascondere
l'Olandese, però, quella cazzo di cosa è così grande che si vede con un
telescopio economico. Girano voci sul ritorno di persone da Paradiso ed
è risaputo che i kristang non sono più al comando, specie ora che stiamo
evacuando Lione e Hangzhou. Vorremmo che tenessi un profilo basso e non
dicessi più del dovuto, ma dovresti vedere la tua famiglia prima di
uscire dai confini del pianeta.»

«È vero, Joe», aggiunse Skippy, «siamo ben oltre il punto in cui
potremmo nascondere quello che è successo, se i kristang dovessero
tornare qui. O tutto o niente.»

\hfill\break

Iniziavo a temere le mattine. Quando mi svegliavo, e con Skippy non era
possibile fingere di dormire, dovevo affrontare qualsiasi problema
avesse avuto durante la notte. Quella mattina era particolarmente
allegro: «Buongiorno. Ehi, ho buone notizie per te. Colonnello Joe, che
tu ci creda o no, gli umani potrebbero non essere solo dei batteri
generici. Penso che la vostra specie abbia inventato una perdita di
tempo che, per quanto ne so, è unica in questa galassia. È
impressionante».

«Facebook? Video di gattini?» Tirai a indovinare. «Solitario al
computer?» Skippy continuava a dire di no. «Ah, andiamo», mi crollarono
le spalle, «il porno devono avercelo anche altre specie.»

«Non è il porno. Sono i fantasport.»

«Mi stai prendendo in giro.''

«Non ti prendo in giro. Nessun'altra specie che conosco spende così
tanto tempo ed energie nello sport, per non fare sport.»

«Ehm.» Non avevo una risposta per questo. I fantasport erano il motivo
di vanto dell'umanità? C'era una sorta di ufficio brevetti galattico,
dove potevamo archiviare la nostra invenzione? Volevo incassare questo
guadagno inaspettato.

«Il fantabaseball, in particolare, è una cosa che può entusiasmare
un'intelligenza artificiale. Così tante statistiche! Così tante
variabili e permutazioni! Non si possono neanche quantificare tutte! Per
quanto ne sa la vostra specie, comunque. È troppo tardi per creare una
squadra di fantafootball, ma non vedo l'ora che inizi la stagione del
baseball. Mi presterai i soldi per iscrivermi?»

Sbattei le palpebre lentamente, cercando di farmi entrare quel concetto
nel cervello. Skippy, che avrebbe potuto introdursi in qualsiasi sistema
informatico bancario del pianeta, rubare un paio di miliardi di dollari
e coprire le sue tracce così bene che nessuno avrebbe mai capito che i
soldi erano scomparsi, voleva prendere in prestito cinquanta dollari da
me?

«Ehm, non ho contanti con me, ma l'esercito mi deve gli arretrati,
quindi, certo, posso trovare un po' di contanti. A quante fantasquadre
vuoi unirti?»

«Tutte.»

«Tutte?»

«Tutte quelle che ci sono online. Ho impostato una sotto-mente qui per
gestire le mie squadre quando saremo partiti con l'Olandese. Perché no?
Sarà grandioso! Cavolo, farò una strage a questo gioco!»

«Non so di\ldots»

«E aspetta solo di vedere il mio tabellone della March Madness. Inoltre,
voglio andare a Las Vegas, baby! Oooh, potrei sbancare al tavolo del
blackjack. E il poker? Non ne parliamo. Quegli idioti non sapranno cosa
li ha colpiti.»

Guardai incredulo il suo coperchio lucido. Avevo creato un mostro:
«Skippy, non si può andare in giro per Las Vegas».

«Perché no? Non mi piacciono le sbronze e le zoccole, ma il gioco
d'azzardo è il mio forte!» Skippy aveva raccolto un po' troppo slang su
internet. «Lo so, lo so, il tuo stupido governo vuole tenermi segreto.
Puoi mettermi in tasca e ti dirò cosa fare. Posso far vibrare i tuoi
timpani a distanza, in modo che nessun altro possa sentirmi. Farò in
modo che ne valga la pena; puoi tenere tutti i soldi, voglio solo
l'azione. E posso procurarti tutte le zoccole che riesci a gestire.»

«Skippy! Non ho bisogno di zoccole!»

«Davvero? Quand'è l'ultima volta che hai inzuppato il biscotto? È stato
un luuungo periodo senza battere chiodo, eh, cowboy? Tutto lavoro e
niente divertimento fa di Joey un ragazzo spento.»

«Niente zoccole!»

«Mmh. Joe, sei riuscito a sorprendermi. Non ti piacciono le ragazze? Non
è questo che il tuo profilo\ldots»

«Oh, così non va.» Come spiegare gli standard sociali umani di
comportamento a un essere che riteneva la moralità una distrazione?
«Senti, Skippy, mi piacciono le ragazze, mi piacciono davvero tanto le
ragazze, mi piacciono le ragazze come persone. Esseri umani. Non ho
niente contro le, ehm, squillo, è solo che non m'interessano. Mi piace
parlare con le ragazze, ok? Non solo buttare soldi sul letto e, ehm,
sai, farlo. E io sono in servizio attivo. Non posso andare in giro per
Las Vegas a rubare soldi, perché è quello che faremmo.»

«Non è vero. Anche per me c'è in gioco un piccolo elemento di
probabilità, è ciò che lo rende una sfida. Perché è ``rubare'' se io
gioco a blackjack, ma non quando il casinò trucca le carte contro i
giocatori?»

Non sapevo cosa rispondergli: «Skippy, ti prometto che chiederò a,
ehm\ldots», ora che ci pensavo, non avevo idea di chi fosse incaricato
di occuparsi di Skippy, a parte la presidentessa in persona. Che non ne
aveva il tempo. Certo, tutte le persone coinvolte avrebbero voluto che
ne fosse incaricato il loro servizio militare o agenzia. «Chiederò che»,
questo aggirava il problema del destinatario della mia richiesta, «tu
faccia qualche viaggetto. Possiamo dire che è, ehm, per questioni di
familiarità culturale o qualcosa del genere.»

«O puoi dire alla tua pres che o vado a Las Vegas o vado in Cina e
sbanco i casinò di Macao, già che ci sono. Oooh, oppure possiamo andarci
insieme! Di' alla gente che è per confrontare culture, o qualche
stronzata del genere. I croupier del blackjack a Macao devono parlare
inglese per te, giusto? Se no, posso insegnarti.»

Potevo vedermi mangiare aspirina come Tic Tac, se avessi dovuto stare
con Skippy ventiquattr'ore su ventiquattro, sette giorni su sette. «Ho
detto: chiederò. Non so se andremo in Cina, potremmo non avere tempo.
Per favore, non fare nulla che potrebbe metterci nei guai. Sai cosa? Se
t'interessa così tanto il calcolo delle probabilità, perché non mi dici
i numeri vincenti del biglietto della lotteria?» Se c'erano ancora
lotterie negli Stati Uniti, molte cose potevano essere cambiate da
quando me n'ero andato.

«Troppo facile.»

«Facile? Sono numeri del tutto casuali! Usano palline da ping-pong.»

«Sì, ovvio, sembra casuale, perché fermarsi\ldots{} ah, continuo a
dimenticare quanto sia lineare il pensiero della tua specie. Non hai
idea di come la quantità\ldots{} Mmh, cazzo. Non posso dirti nulla senza
interferire in modo drastico con lo sviluppo della vostra specie.»

«Ancora restrizioni nella tua programmazione?»

«No, è immorale.» Il suo tono di voce implicava un ``ma tu pensa'' che
non espresse ad alta voce.

«Immorale? Detto da te?»

«So che per te è difficile da credere, ma, riguardo alle questioni
importanti, sono molto severo sulla moralità.»

«Tipo, giocare a poker e derubare i casinò?»

«È moralmente sbagliato lasciare che gli idioti si tengano i loro soldi,
altrimenti, non imparano nulla. E i casinò? Andiamo, sono loro a
derubare la gente. Ho detto le questioni morali importanti, no? I soldi
non sono importanti.»

\hfill\break

Mentre andavo a un incontro con gli ambasciatori di Gran Bretagna e
Cina, io e Skippy fummo avvicinati in un corridoio del centro conferenze
dal dottor Constantine.

«Sergente! Sergente Bishop!», chiamò, senza fiato. «Ho bisogno di
parlare con lei!»

«Oh, merda», mormorai sottovoce, «non sapevo che questo idiota fosse
qui.»

«Io lo sapevo», mormorò Skippy, «ho riprogrammato la sua sveglia perché
facesse tardi, ma si è alzato comunque in tempo, cazzo.»

«Sergente, volevo dirle che non vedo l'ora di partire per la missione.
Spero di avere l'opportunità di discutere ancora con lo\ldots»,
Constantine inciampò sulle sue parole, «con lo\ldots{} lo Skippy.»

Gli rivolsi il mio miglior sguardo di ghiaccio: «Sono il colonnello
Bishop», indicai le aquile d'argento sul mio colletto. Aquile che erano
di nuovo rivolte verso il ramo d'ulivo, invece che verso le frecce.
«Sono al comando della missione. E non ho visto il suo nome sulla lista
dei volontari.»

«Cosa?», disse Constantine, scioccato per il fatto che fossi io al
comando o che il suo nome non fosse sulla lista, o entrambi. «Le
assicuro\ldots»

«Il suo nome era sulla lista», spiegò Skippy, «ma io l'ho cancellato dal
database. Al colonnello Joe non piaci, a me non piaci, quindi non
verrai.»

«Non può essere», sputò Constantine. «Serg\ldots{} colonnello Bishop»,
disse il mio grado come se non ci potesse credere, «certo capisce che
una\ldots{} una persona del suo rango», stava davvero lottando con tutta
la questione delle competenze sociali, «non può permettere che i
sentimenti personali interferiscano con la necessità di avere le persone
più qualificate sulla\ldots{} a bordo della nave, sotto\ldots{} sotto il
suo comando. Senza offesa, ma parlerò con i suoi superiori, capiranno
che è vitale che abbiamo i nostri migliori uomini in questo viaggio.» Se
intendeva davvero quello che aveva detto, aveva molto bisogno che gli
aprissero gli occhi sul concetto di ``senza offesa''.

«Parla con chi ti pare», disse acido Skippy. «L'unico modo per arrivare
all'Olandese è su una navicella thuranin che piloterò io e, se ci sarai
tu a bordo, quella navicella non andrà da nessuna parte.»

«Dottor Constantine, ho fatto una ricerca su di lei dopo che ci siamo
incontrati a Colorado Springs», ammisi, «e, da quello che ho letto, lei
è una delle menti più brillanti del XXI secolo», riconobbi. Il tizio
aveva iniziato il college all'Istituto di tecnologia del Massachussets
quando aveva solo quattordici anni, e aveva già vinto diversi premi
scientifici importanti prima di allora. Non riuscivo nemmeno a capire i
titoli degli articoli che aveva scritto, figuriamoci afferrarne il
contenuto. Il suo volto s'illuminò di un sorriso prima che potessi
schiacciare le sue speranze. «Purtroppo i suoi colleghi pensano che lei
sia anche uno dei più grandi stronzi del XXI secolo. Dà sui nervi a
tutti quelli con cui lavora.» Le migliori istituzioni scientifiche del
pianeta lo avevano licenziato, o costretto, invitato, a lasciarle. In un
campo che doveva generare tanti ego titanici quante brillanti scoperte
scientifiche, quanto dovevi essere stronzo per distinguerti al punto che
persone di enorme talento non riuscissero a lavorare con te? «Ha ragione
sul fatto che ho bisogno dell'equipaggio più qualificato, sotto il mio
comando», enfatizzai l'ultima parte. «Sarò responsabile di settanta
persone, in una pericolosa missione lontana dalla Terra, per due anni e
mezzo.» Due anni e mezzo, se saremmo stati fortunati. «Le persone in
questa missione non solo devono essere tra le migliori nel loro campo,
ma devono anche andare d'accordo le une con le altre, a stretto
contatto. Questo la esclude dalla lista.»

«L'altro tuo problema è», aggiunse Skippy in tono allegro, «che la tua
unica qualifica per la missione è il fatto che tu sia, beh, un po' più
intelligente della scimmia media. Il che, al mio cospetto, non vale
niente. Dottore, tu hai una considerazione troppo alta di te stesso.
L'unica differenza tra te e Joe è che Joe è come il cane che guarda
attraverso il parabrezza, e tu sei il cane che sporge la testa fuori dal
finestrino. Avrai una visuale di poco migliore, ma non la capirai
meglio.»

«Non consiglierei comunque di sporgere la testa fuori dal finestrino»,
aggiunsi, «non c'è brezza nello spazio.»

«Joe non ha tutti i torti», concluse Skippy. «E poi, se ci fosse brezza,
sbaveresti su tutta la nave dietro di te.»

Constantine mi guardava, come se si aspettasse che la solidarietà
interspecifica mi avrebbe fatto perorare la sua causa. Quello che non
aveva capito è che, di fatto, il concetto universale era che un idiota è
un idiota, a prescindere dalla specie.

«Farò in modo di inviarle una cartolina», mi offrii. «Ci sarà scritto:
``Ci stiamo divertendo un mondo, lieto che lei non sia qui''.»

\hfill\break

L'incontro con l'ambasciatore cinese iniziò quasi in contemporanea col
mio incontro in corridoio con il dottor McPippas. Sembrava che, mentre i
governi coinvolti avevano concordato sull'invio dell'Olandese nello
spazio, il che a mio avviso equivaleva a un eureka meritevole di un ``ma
tu pensa'', ci fosse meno accordo sull'opportunità di mettere me al
comando della missione. L'ambasciatore britannico mi strinse la mano, mi
augurò buona fortuna e ci tenne a menzionare la squadra Sas, che era il
contributo del suo paese all'equipaggio militare dell'Olandese. Se il
governo britannico aveva delle riserve sul fatto che fossi io a
comandare la missione, erano tutti troppo educati per dirlo. O stimavano
che non valesse la pena litigare.

L'ambasciatore cinese fu più diretto. La Cina aveva un generale
dell'Aeronautica a due stelle, un tizio con un curriculum
impressionante, che pensavano dovesse comandare la missione. I cinesi
non avevano nulla in contrario al fatto che salissi a bordo, specie
nelle vesti di collegamento con Skippy, ma, anche se ora ero di nuovo un
colonnello, un generale a due stelle mi avrebbe superato per grado, con
una notevole quantità di autorità.

Gerald Schmidt era un consigliere speciale della Casa Bianca, che era
stato nominato per facilitare le relazioni con i nostri alleati rispetto
al viaggio dell'Olandese. Schmidt cercò di negoziare con i cinesi, ma il
generale Brenner non ne voleva sapere: «Allora promuoveremo il
colonnello Bishop a generale a tre stelle». Lanciò un'occhiataccia. «O,
diamine, a cinque stelle, se necessario. Non giocheremo a questo gioco
con voi.»

L'ambasciatore cinese doveva avere deciso che il tempo per la diplomazia
era finito. «L'arroganza americana è sorprendente. Il vostro paese non è
una superpotenza nello spazio, eppure vi comportate come se\ldots»

«Ehi, ehi», disse Skippy quasi gridando, «smettetela di blaterare!» Era
diventata una delle sue espressioni preferite. «Voi scimmie senza
cervello fate quello che volete coi titoli e le uniformi, per me non
significano un bel niente. Il grado di Joe potrà anche essere quello di
grande pallone gonfiato oppure di Bobo il Pagliaccio, e sarà comunque il
capitano della nave, capito? Se voi scimmie pulciose volete venire,
allora accettate che il colonnello Joe sia al comando.»

Bobo il Pagliaccio? Mi chiesi in quel momento -- era stupefacente quanto
la mia mente potesse vagare --: i pagliacci hanno rango? Un pagliaccio
con un grosso naso rosso supera forse un pagliaccio con un\ldots{}

«È giovane e inesperto», i cinesi parlavano piano, in modo diplomatico,
rivolgendosi direttamente a Skippy. «Cos'ha di speciale il colonnello
Bishop?»

Mi stavo chiedendo la stessa cosa.

«Non ho bisogno di giustificarmi con te», Skippy storse il naso, «ma
continueresti a rompermi le scatole in proposito, quindi te lo dirò. Joe
è l'unico della tua specie sottosviluppata che mi ha trattato come un
essere senziente, come una persona, fin dall'inizio. Voialtre scimmie mi
considerate una macchina e avete paura di me. Joe mi ha dato un nome.
Prima ho sempre avuto solo una designazione. Ora ho un nome. Joe, tu hai
detto che sono uno stronzo e il motivo per cui l'hai detto è che mi
consideri al livello di un essere del tutto senziente. Una macchina non
può essere stronza, solo una persona può esserlo. Mi tratti come una
persona, come un tuo pari.»

Ero davvero commosso: «Grazie, Skippy, non mi hai mai, ehm, detto niente
al riguardo prima».

«Speravo che l'avresti capito da solo. Ahimè, non sarebbe mai successo
con il tuo cervello lento.»

«Eh, sei ancora uno stronzo, vedo.»

«Dimostrando con esattezza il mio punto di vista», disse Skippy
compiaciuto.

«Signor Skippy», l'ambasciatore cinese aveva un'espressione dolente sul
volto mentre chiamava Skippy un'intelligenza artificiale
super-brillante. Cosa che mi fece riflettere. ``Skippy'' non significava
nulla in cinese, quindi perché gli importava? A meno che non fossero
imbarazzati per noi. Chissà. «È comprensibile che lei sia più a suo agio
avendo una persona familiare a bordo della nave, ma è la persona
migliore per comandare? Entrare in contatto con il Collettivo è la sua
priorità, non dovrebbe avere un comandante che possa assicurare il
successo della missione?»

«Certo, non fa una piega», concordò Skippy. «Dammi una lista di
candidati che hanno più esperienza di Joe al comando di astronavi aliene
e la esaminerò. Fino ad allora, chiudi quella cazzo di bocca.»

Era chiaro che l'ambasciatore cinese non avrebbe seguito il consiglio di
Skippy di tenere la bocca chiusa, così alzai un dito indice. «Signor
ambasciatore, un momento, per favore?» Feci segno agli altri americani
di riunirsi nell'angolo della stanza. «E se io comandassi la nave e il
tenente Chang fosse a capo dell'equipaggio?»

Schmidt inclinò la testa verso il generale Brenner: «Questo
permetterebbe ai cinesi di salvare la faccia. Cosa ne pensi?».

«Penso che non abbiamo bisogno dei cinesi in questa missione, e se non
vogliono collaborare possono restare a casa», ringhiò Brenner, guardando
la sua controparte cinese all'estremità opposta della stanza.

«La presidentessa vuole che questa sia una missione multinazionale ed è
nostro compito accontentarla, con successo», ricordò Schmidt in tono
gentile. «Dobbiamo lavorare insieme per ricostruire questo pianeta e,
una volta che noi umani saliremo sulle nostre navi lassù», indicò il
soffitto, «avremo bisogno di una forza umana unificata. Vi sta bene
questo accordo?»

«Cinesi al comando degli americani?», schernì Brenner.

«Con me al comando generale, signore», feci notare. «E nulla avviene
senza che Skippy sia d'accordo, in ogni caso.» La mascella di Brenner
faceva avanti e indietro, come se stesse masticando qualcosa che non
riusciva a ingoiare. «Conosco Chang, signore, è un bravo ragazzo.» Un
bravo ragazzo? Tutto quello che avevo fatto io lì era ricordare a
Brenner che stava consegnando il comando della missione a un giovane
sergente inesperto, a prescindere da quali mostrine ci fossero sulla mia
uniforme.

«Colonnello», Brenner mi fissò negli occhi, ed ero determinato a non
tirarmi indietro, «dopo che sarai saltato via, sarai da solo. Penso che
fare in modo che i nostri alleati salvino la faccia sia meno importante
che mantenere una chiara catena di comando, e penso che questo accordo
stia creando problemi, ma la decisione spetta a te. Dopo il salto, ogni
decisione spetterà a te.» Fu difficile non distogliere lo sguardo quando
lo disse, ma non lo feci.

Merda. Ora stavo dubitando di me stesso. Era troppo tardi per cambiare
idea: «Posso fare in modo che funzioni, signore. Se Chang crea problemi,
chiederò a Skippy di chiuderlo nella sua cabina». Chang sapeva, per
esperienza, che non era successo nulla a bordo dell'Olandese senza che
io e Skippy lo sapessimo e lo approvassimo. Un altro ufficiale avrebbe
potuto mettersi in testa di prendere il controllo. Chang non lo avrebbe
fatto.

Inoltre, non l'avevo detto a nessuno, ma avere Chang che si occupava
dell'equipaggio mi sollevava da tutte le stronzate amministrative
associate. Un incentivo, per quanto mi riguardava.

«Siamo d'accordo, allora», Schmidt annuì rivolto a me. «È stata una
buona idea, colonnello Bishop. Avrà bisogno di abilità diplomatiche per
guidare l'equipaggio internazionale là fuori. Forse lei è la persona
giusta per l'incarico.»

«Grazie, signore.»

Schmidt aggrottò le sopracciglia: «Visto che abbiamo un momento, ho
bisogno di chiedere del dottor Constantine».

Brontolai dentro di me, pronto a combattere: «La sua personalità lo
rende inadatto per un lungo viaggio in ambienti chiusi, signore».

Con mia sorpresa, Schmidt sorrise: «Bene, va bene. Ok, è fuori dalla
lista».

«A posto così?», chiesi.

«Sì. Constantine ha dei sostenitori potenti, ma se il comandante della
missione dice che il suo profilo personale non è adatto, non credo che
qualcuno possa discutere su questo. La Casa Bianca stava cercando una
scusa per lasciarlo fuori dalla lista della missione.»

L'ambasciatore cinese accettò, con entusiasmo, l'idea di mettere Chang a
capo dell'equipaggio, e propose che fosse nominato in via ufficiale
vicecomandante, per formalizzare la sua autorità, e io accettai. Pace e
armonia tra l'equipaggio, il dottor Constantine fuori dalla lista,
un'opportunità per vedere i miei genitori prima che partissimo. Erano
tutti buoni presupposti per il nostro viaggio. Un viaggio dalla durata
indeterminata, in una galassia ostile, a bordo di una nave rubata, con
pezzi di ricambio limitati. Una nave che non capivamo, guidata da
un'intelligenza artificiale aliena distratta, che aveva solo una vaga
idea di dove fossimo diretti.

Cosa sarebbe potuto andare storto?